%%%%%%%%%%%%%%%%%%%%%%%%%%%%%%%%%%%%%%%%%%%%%%%%%%%%%%%%%%%%%%%%%%%%%%%%%%%%%%%
%  正標数 log parabolic geometry における有限レベル現象
%  作業ドラフト  ver.0.4
%
%  コンパイル: lualatex (ltjsarticle)
%  ※ upLaTeX を使う場合は documentclass を
%     \documentclass[uplatex,dvipdfmx,a4paper,11pt]{jsarticle}
%     に差し替える。
%
%  変更履歴
%  --------
%  ver.0.4 (2026-09-06):
%    * \S6 の数え上げを Wakabayashi の集合 {}^\dagger C_N ベースに全面改訂。
%    * \S6 の N=2 の数値を訂正（前稿 ver.0.3 の 610 / 2262 は誤りかつ未検証。
%      正しい値 225 / 1666 等を computations/count_dagger_C.py で再検証）。
%    * \S8 を「持ち上げの障害は行列式に尽きる」に再編。
%      thm:M03, cor:general を「仮定 (W10.3)」つきの条件付き定理に変更
%      （Wakabayashi Prop. 10.3.3 未確認 = L2 の穴を明示）。
%      thm:det, cor:pi1（障害の det 局在）を追加。
%    * 付録 A（訂正の履歴）、付録 B（文献確認事項）を追加。
%    * ステータスマクロ \GAP, \NeedsCheck を追加。
%  詳細な地位区分は下の「主張の地位について」を参照。
%%%%%%%%%%%%%%%%%%%%%%%%%%%%%%%%%%%%%%%%%%%%%%%%%%%%%%%%%%%%%%%%%%%%%%%%%%%%%%%

\documentclass[a4paper,11pt]{ltjsarticle}

\usepackage{amsmath,amssymb,amsthm,mathrsfs}
\usepackage{booktabs}
\usepackage{array}
\usepackage{enumitem}
\usepackage{xcolor}
\usepackage[hidelinks]{hyperref}

%%%%%%%%%%%%%%%%%%%%%%%%%%%%%%%% 定理環境 %%%%%%%%%%%%%%%%%%%%%%%%%%%%%%%%%%%%%%

\theoremstyle{plain}
\newtheorem{theorem}{定理}[section]
\newtheorem{proposition}[theorem]{命題}
\newtheorem{lemma}[theorem]{補題}
\newtheorem{corollary}[theorem]{系}
\newtheorem{conjecture}[theorem]{予想}

\theoremstyle{definition}
\newtheorem{definition}[theorem]{定義}
\newtheorem{example}[theorem]{例}
\newtheorem{remark}[theorem]{注意}
\newtheorem{question}[theorem]{問題}

% ステータスタグ（清書時に削除する）
%   \Todo       : 今後行う作業一般
%   \GAP        : 数学的な証明の穴（仮定を追加しても閉じないもの）
%   \NeedsCheck : 文献・計算・定義等の確認待ち
%   \Unver      : 現時点で未検証の主張
%   \Known      : 既知または確認済み事項
\newcommand{\Known}[1]{\marginpar{\footnotesize\textsf{[既知: #1]}}}
\newcommand{\Todo}[1]{\textcolor{red}{\textsf{[TODO: #1]}}}
\newcommand{\Unver}{\textsf{[未検証]}}
\newcommand{\GAP}[1]{\textcolor{orange!85!black}{\textsf{[GAP: #1]}}}
\newcommand{\NeedsCheck}[1]{\textcolor{blue!70!black}{\textsf{[要確認: #1]}}}

%%%%%%%%%%%%%%%%%%%%%%%%%%%%%%%% 記号 %%%%%%%%%%%%%%%%%%%%%%%%%%%%%%%%%%%%%%%%%%

\newcommand{\bbA}{\mathbb{A}}
\newcommand{\bbP}{\mathbb{P}}
\newcommand{\bbZ}{\mathbb{Z}}
\newcommand{\bbF}{\mathbb{F}}
\newcommand{\cD}{\mathcal{D}}
\newcommand{\cE}{\mathcal{E}}
\newcommand{\cF}{\mathcal{F}}
\newcommand{\cL}{\mathcal{L}}
\newcommand{\cM}{\mathcal{M}}
\newcommand{\cN}{\mathcal{N}}
\newcommand{\cO}{\mathcal{O}}
\newcommand{\cT}{\mathcal{T}}
\newcommand{\sX}{\mathscr{X}}
\newcommand{\sP}{\mathscr{P}}
\newcommand{\fg}{\mathfrak{g}}
\newcommand{\fp}{\mathfrak{p}}
\newcommand{\fS}{\mathfrak{S}}
\newcommand{\Res}{\operatorname{Res}}
\newcommand{\Pic}{\operatorname{Pic}}
\newcommand{\End}{\operatorname{End}}
\newcommand{\Hom}{\operatorname{Hom}}
\newcommand{\Coker}{\operatorname{Coker}}
\newcommand{\rk}{\operatorname{rk}}
\newcommand{\Op}{\mathcal{O}p}
\newcommand{\Zzz}{\mathrm{Zzz\dots}}
\newcommand{\dlog}{d\!\log}
\newcommand{\Dlog}[1]{\cD^{(#1)}_{X^{\log}}}
\newcommand{\Dbr}[1]{\breve{\cD}^{(#1)}_{R}}

%%%%%%%%%%%%%%%%%%%%%%%%%%%%%%%%%%%%%%%%%%%%%%%%%%%%%%%%%%%%%%%%%%%%%%%%%%%%%%%

\title{正標数における対数的レベル理論と\\有限レベル現象および持ち上げ障害の局在について\\
\large ---- 作業ドラフト ver.\,0.4 ----}
\author{}
\date{\today}

\begin{document}
\maketitle

\begin{quote}\small\sffamily
\textbf{主張の地位について.}\ 本稿は作業ドラフトであり、主張の地位を次のように区別する。
\emph{[文献]}：外部文献の結果（証明は追っていない。原論文 PDF を直接確認していない
ものは \NeedsCheck{} で示す）。
\emph{[本研究・定理]}：本稿で完全な証明を与えたもの。
\emph{[本研究・条件付き定理]}：明示した仮定のもとで完全な証明を与えたもの。
\emph{[本研究・部分証明]}：証明の骨格はあるが穴を自覚しているもの（\GAP{} で示す）。
\emph{[予想]}：証拠はあるが証明がないもの。
数値は計算機で実際に実行した結果のみを確定値として記す（付録 \ref{app:num}）。
\end{quote}

\begin{abstract}
$k$ を標数 $p>0$ の代数閉体、$(X,D)$ を $k$ 上の log smooth pair とする。
Berthelot--Montagnon の有限レベル対数微分作用素環 $\cD^{(m)}_{X^{\log}}$ に対する
dormant 加群（$p^{m+1}$-曲率が消える加群）について、
「有限レベルでは成立するが無限レベルでは成立しない」現象、
すなわち与えられた dormant 対象がどこまで高いレベルに持ち上がるかという問題を扱う。
主結果は次の三点である。
(1) 対数的な場合、dormant 対象の局所不変量 exponent は $\bbZ/p^{N}\bbZ$ に値をとり、
レベルの上昇は exponent の $p$ 進展開に桁を加えることに他ならない（\S\ref{sec:local}）。
したがって持ち上げの障害は純粋に大域的であり、
その障害は行列式束の次数に関する合同式として書ける（命題 \ref{prop:degree}）。
(2) 3点付き射影直線 $\sP=(\bbP^1;0,1,\infty)$ 上では
dormant $\mathrm{PGL}_2$-oper の個数はレベル $N=1$ で $p(p^2-1)/24$ であり
（命題 \ref{prop:count}）、
かつ Wakabayashi の径数付けとレベル降下の両立性（仮定 (W10.3)、本稿では未確認）
のもとで制限写像 ${}^{N+1}\Op\to{}^{N}\Op$ は全射である（定理 \ref{thm:M03}）。
すなわちこの仮定のもとで $\sP$ 上には反例が存在しない。
(3) 同じ仮定のもと、totally degenerate な曲線上でも制限写像は全射であり
（系 \ref{cor:general}）、その範囲では $(g,r)$ によらず
「ちょうどレベル $N$」の $\mathrm{PGL}_2$-oper 反例は存在しない。
持ち上げの障害は行列式束に局在すると考えられ
（定理 \ref{thm:det}、ただし $p\mid n$ の場合は未解決）、
$\pi_1(G)$ が $p$ 群を含まなければ $\Pic(X)/p^N$ の障害に帰着する（系 \ref{cor:pi1}）。
一般の smooth 曲線での全射性は未解決である（\S\ref{sec:open}）。
$p=5$ における dormant $\mathrm{PGL}_2$-oper 5 個の完全な明示リストを与える（\S\ref{sec:p5}）。
\end{abstract}

\tableofcontents

%%%%%%%%%%%%%%%%%%%%%%%%%%%%%%%%%%%%%%%%%%%%%%%%%%%%%%%%%%%%%%%%%%%%%%%%%%%%%%%
\section{序}\label{sec:intro}

\subsection{動機}

$k$ を標数 $p>0$ の代数閉体、$X$ を $k$ 上の滑らかな多様体とする。
$X$ 上の平坦接続 $(\cE,\nabla)$ に対し、
その $p$-曲率 $\psi_{(\cE,\nabla)}$ が消えることは
Cartier 降下により $(\cE,\nabla)\cong F^{*}(\cE')$（標準接続つき）と同値である。
さらに $\cE'$ 上の接続を選び直して降下を繰り返すことができるかを問うと、
Berthelot \cite{Ber1} の有限レベル微分作用素環の増大列
\[
  \cD^{(0)}_X\to\cD^{(1)}_X\to\cdots\to\cD^{(m)}_X\to\cdots,
  \qquad \varinjlim_m\cD^{(m)}_X=\cD^{(\infty)}_X
\]
が現れ、$\cD^{(\infty)}_X$-加群（stratified sheaf）が無限レベルの対象となる。

非対数的な場合、この階層は Frobenius 降下の階層と完全に一致する：
$\cD^{(m)}_X$-加群で $p^{m+1}$-曲率が消えるもの（dormant なもの）の圏は
$\cO_{X^{(m+1)}}$-加群の圏と同値である。
したがって「レベル $N$ まで持ち上がる」とは
「$F^{(N)*}$ の像に入る」ということに尽きる。

本稿の出発点は、この一致が\textbf{対数的な場合には崩れる}という観察である。
$X^{\log}=(X,D)$ を log smooth pair とすると、
対数接ベクトル場は toral である：局所座標で $\theta=x\partial_x$ とすると
\begin{equation}\label{eq:toral}
  \theta^{[p]}=\theta .
\end{equation}
実際 $\theta(x^n)=nx^n$ より $\theta^p(x^n)=n^px^n=nx^n$（Fermat の小定理）。
この一点が以下すべての特異性の源である。

\subsection{出発点となる計算}\label{ss:starting}

$X=\bbA^1$、$D=\{0\}$、$\cE=\cO_X$、$\nabla=d+a\,\dlog x$（$a\in k$）とする。
$\nabla_\theta=\theta+a$ であり、$\theta$ とスカラー $a$ は可換なので
可換環における Frobenius の加法性から
\begin{equation}\label{eq:AS}
  \psi(\theta)=(\theta+a)^p-(\theta+a)
  =\theta^p+a^p-\theta-a
  =(\theta^{[p]}-\theta)+a^p-a
  =a^p-a .
\end{equation}
よって
\[
  \psi=0\iff a^p=a\iff a\in\bbF_p .
\]
対照的に、非対数的な $\bbA^1$ 上で $\nabla=d+a\,dx$ とすると
$\partial^{[p]}=0$ かつ $\partial^p=0$（$\cO$ 上）だから $\psi(\partial)=a^p$ であり、
消えるのは $a=0$ のときのみである。
Artin--Schreier 方程式 $a^p-a=0$ が現れるのは
toral 性 \eqref{eq:toral} の直接の帰結であり、
対数構造が \emph{ちょうど $\bbF_p$ 個の} 非自明な dormant 対象を生むことを意味する。

一方、log Frobenius $F$ はモノイド上で $m\mapsto m^p$ なので
\[
  F^*(\dlog t)=\dlog(x^p)=p\,\dlog x=0 ,
\]
すなわち $F^*\Omega^1_{X'}(\log D')\to\Omega^1_X(\log D)$ は零写像である。
したがって素朴な Frobenius 引き戻しは留数 $0$ の対象しか生まない。
$a\neq0$ なる \eqref{eq:AS} の対象は $p$-曲率が消えるにもかかわらず
素朴な $F^*$ の像には入らない。ここに「有限レベル／無限レベル」の分岐の芽がある。

\begin{remark}[当初の解釈の訂正]\label{rem:early-misread}
本研究の初期には、この観察を「対数的な場合には『$\cD^{(1)}$-作用の存在』と
『$p$-曲率 $0$』が乖離する」と解釈していた。これは誤りであり、
素朴な $F^*$ と $\cD^{(1)}$-加群構造への持ち上げの混同であった。
実際には $a\in\bbF_p^\times$ の対象も $\cD^{(1)}$-加群として持ち上がる
（注意 \ref{rem:binom}）。すなわち上の非包含は「素朴な $F^*$ の像に入らない」
というだけの主張であって、レベル上昇の障害ではない。
訂正の経緯は付録 \ref{app:errata} を参照。
\end{remark}

\subsection{本稿の結論の要約}

上の観察を精密化すると、次の描像が得られる（\S\ref{sec:local}）。
\begin{itemize}[leftmargin=1.6em]
\item dormant な対数対象の局所不変量は
      \textbf{exponent} $d\in\bbZ/p^{N}\bbZ$ であり、
      その $p$ 進展開 $\tilde d=\sum_{l}p^l\tilde d_l$ の第 $l$ 桁が
      $l$ 段目に降下した対象の留数（モノドロミー）を与える。
\item レベルを $N$ から $N+1$ に上げることは、
      exponent を $\bbZ/p^{N}\to\bbZ/p^{N+1}$ に持ち上げることに他ならない。
      これは\textbf{局所的には常に可能}（各分岐点で $p$ 通りの選択）。
\item よって障害は大域的であり、
      行列式束の次数に関する合同式（命題 \ref{prop:degree}）に集約される。
      marked point が $r\ge1$ 個あれば持ち上げの自由度が $p^r$ 通りあるので
      合同式は解けることが多く、$r=0$ のときにのみ $\Pic$ の障害が残る。
\end{itemize}
これは当初の予想と逆であり、
\textbf{対数構造は障害を消す方向に働く}というのが本稿の主張である。

\subsection{規約に関する注意}\label{ss:convention}

レベルの添字には二つの流儀があり、混乱の元なので明記する。
\begin{itemize}[leftmargin=1.6em]
\item \textbf{Berthelot 流}：$\cD^{(m)}$ の $m$。
\item \textbf{Wakabayashi 流 \cite{Wak}}：「レベル $N$ の oper」$=$ $\cD^{(N-1)}$-加群構造つき oper、
      「dormant」$=$ $p^N$-曲率 $\nabla_{\langle p^N\rangle}$ の消滅。
\end{itemize}
非対数的な場合、
$\mathfrak{Mod}(\cD^{(m)}_S)\simeq\mathfrak{Mod}(\cD^{(m-l)}_{S^{(l)}})$ かつ
$p^{m+1}$-曲率が $p^{m-l+1}$-曲率に対応する \cite[Prop.~2.1.5]{Wak} ので
\[
  \text{dormant レベル }N \iff F^{(N)*}(\cO\text{-加群})
\]
となり、Wakabayashi 流の「dormant レベル $N$」は
「$N$ 回 Frobenius 降下できる」という素朴な意味と一致する。
\textbf{本稿では Wakabayashi 流に統一する。}
すなわち「レベル $N$」は $\cD^{(N-1)}$-作用を指し、
「dormant」は $p^N$-曲率の消滅を指す。

\begin{remark}[規約の照合状況]\label{rem:conv-status}
「レベル $N$ $=$ $\cD^{(N-1)}$」という対応は Wakabayashi \cite[\S4.2]{Wak}
の記述と照合した（2026-09-03）。ただしこの照合は ar5iv 経由の二次的なものであり、
原論文 PDF を直接目視したものではない
（\texttt{handoff/app-session-2026-09-03.md} \S0.3）。
Berthelot \cite{Ber1} の原論文における $\cD^{(m)}$ の定義そのもの
（$\partial^{[j]}$ が $\cD^{(m)}$ に属する $j$ の範囲が $j<p^{m+1}$ か否か）、
および正標数 $\bbZ_p$ 値留数・logarithmic holonomy の各研究で用いられる規約
との関係は、本稿では確定していない。
詳細は \texttt{notes/topics/00-common/conventions.tex}（C-1, C-1a, C-2）。
本稿は他テーマの規約をここで採用・統一しない。
\NeedsCheck{\cite[\S1.1--1.2]{Ber1} の $\cD^{(m)}$ の定義（本稿では未確認。
付録 \ref{app:litcheck}）。}
\end{remark}

%%%%%%%%%%%%%%%%%%%%%%%%%%%%%%%%%%%%%%%%%%%%%%%%%%%%%%%%%%%%%%%%%%%%%%%%%%%%%%%
\section{記号と準備}\label{sec:notation}

\subsection{log smooth pair と対数微分作用素}

$k=\bar k$、$\operatorname{char}k=p>0$。
$(X,D)$ を log smooth pair とし、局所的に $D=\{x_1\cdots x_r=0\}$ と書く。
$T_X(-\log D)$ の局所基底を
\[
  \theta_i=x_i\partial_i\ (i\le r),\qquad \partial_j\ (j>r)
\]
とすると、制限 Lie 代数構造は
\[
  \theta_i^{[p]}=\theta_i\ \text{(toral)},\qquad \partial_j^{[p]}=0\ \text{(nilpotent)} .
\]

$X^{\log}$ を対応する fs log scheme とし、
$\Dlog{m}$ で Montagnon \cite{Mon} によるレベル $m$ の対数微分作用素環を表す。
$\cD^{(m)}$-加群構造は、$\nabla_{\langle j\rangle}:=\nabla(\partial_{\langle j\rangle})$
（$j\in\bbZ_{\ge0}$）の族で記述される \cite[\S1.3]{Wak}。
$\Dlog{0}$-加群構造を与えることは対数接続を与えることと同値である。

\begin{definition}[$p^N$-曲率、dormant]
$\Dlog{N-1}$-加群 $(\cE,\nabla)$ に対し
${}^p\psi_{(\cE,\nabla)}:=\nabla\circ{}^p\psi_{X^{\log}}$
を $p^N$-曲率とよび、これが消えるとき $(\cE,\nabla)$ は \textbf{dormant} であるという。
$N=1$ のとき、これは対応する対数接続の $p$-曲率の消滅と一致する。
\end{definition}

\subsection{局所モデル}

$R$ を離散付値環、$t$ を素元、$K=\operatorname{Frac}(R)$、$k$ を剰余体とする。
Wakabayashi \cite[\S2.1]{Wak} に従い、$R$ 上の
extended $m$-log derivation を
\begin{equation}\label{eq:logderiv}
  \breve\partial_{\langle j\rangle}(t^n):=q_j!\cdot\binom{n}{j}\cdot t^n,
  \qquad q_j:=\lfloor j/p^m\rfloor
\end{equation}
で定め、対応する環を $\Dbr{m}$ と書く。
$X^{\log}$ の marked point $\sigma$ における $\Dlog{N-1}$ の局所化は
$\breve\cD^{(N-1)}_{\cO_{X,\sigma}}$ と同一視される \cite[Lemma 2.3.3]{Mon}。

\begin{remark}\label{rem:binom}
$\bbZ_p$ における恒等式
$\prod_{j=0}^{p-1}(\theta-j)=p!\binom{\theta}{p}$
の左辺は $\bmod\,p$ で $\theta^p-\theta$ である。
すなわち $\theta^p-\theta^{[p]}$ は $\cD^{(1)}$ の基底元 $\binom{\theta}{p}$ の
$p!$ 倍として持ち上がる。
これがレベル上昇と $p$-曲率の関係の環論的な出所である。
\end{remark}

%%%%%%%%%%%%%%%%%%%%%%%%%%%%%%%%%%%%%%%%%%%%%%%%%%%%%%%%%%%%%%%%%%%%%%%%%%%%%%%
\section{留数の強制的半単純性}\label{sec:residue}

まずレベル $1$（通常の対数接続）における基本的な制約を述べる。

\begin{lemma}[留数の toral 性]\label{lem:toral}
$(\cE,\nabla)$ を $(X,D)$ 上の対数接続、
$R_i=\Res_{D_i}(\nabla)\in\End(\cE|_{D_i})$ とする。
$p$-曲率 $\psi$ が消えるならば $R_i^{\,p}=R_i$ である。
特に $R_i$ は半単純であり、その固有値は $\bbF_p$ に属する。
\end{lemma}

\begin{proof}
$\nabla_{\theta_i}(x_ie)=x_i(\nabla_{\theta_i}+1)e$ より
$\nabla_{\theta_i}$ は $x_i\cE$ を保つ。
よって $\cE/x_i\cE$ 上で $\nabla_{\theta_i}^{\,p}$ は $R_i^{\,p}$ を誘導する。
$\psi(\theta_i)=\nabla_{\theta_i}^{\,p}-\nabla_{\theta_i}$ を $D_i$ に制限すれば
$R_i^{\,p}-R_i=0$。
$T^p-T$ は分離多項式なので $R_i$ は対角化可能で、固有値は $\bbF_p$ に属する。
\end{proof}

\begin{corollary}\label{cor:nilpotent}
$\psi=0$ かつ $R_i$ が冪単（すなわち $R_i-\mathrm{id}$ が冪零）ならば $R_i=0$。
特に、非自明な冪零留数をもつ対数接続の $p$-曲率は消えない。
\end{corollary}

\begin{proof}
冪零元 $N$ に対し $N^p=0$ なので $R^p=R$ から $R=0$。
\end{proof}

\begin{example}
$\nabla_\theta=\theta+N$、$N$ を定数冪零行列とすると
$\psi(\theta)=N^p-N=-N\neq0$。
\end{example}

\begin{remark}[主束版]
主 $G$ 束の場合、$\Res\in\fg_{E}$ が $\Res^{[p]}=\Res$ を満たす、
すなわち $\Res$ は toral 元であり、これは準同型 $\mu_p\to G$ を与える。
共役類として
\[
  \Res\in\mathfrak{c}(\bbF_p)=\bigl(X_*(T)\otimes\bbF_p\bigr)/W .
\]
$G=\mathrm{PGL}_2$ では $\mathfrak{c}(\bbF_p)=\bbF_p/\{\pm1\}$ であり、
これが \S\ref{sec:opers} の radius の概念的な正体である。
\end{remark}

\begin{remark}[注意点への回答]
系 \ref{cor:nilpotent} により、
「留数が半単純か冪単かで挙動が分かれる」という当初の予想は、
dormant な世界では\textbf{半単純の場合しか存在しない}という形で決着する。
冪単留数の対象は $p$-曲率が非零な世界（Mochizuki の nilpotent indigenous bundle）
に属し、本稿の理論とは交わらない。
\end{remark}

%%%%%%%%%%%%%%%%%%%%%%%%%%%%%%%%%%%%%%%%%%%%%%%%%%%%%%%%%%%%%%%%%%%%%%%%%%%%%%%
\section{局所理論：exponent と $p$ 進展開}\label{sec:local}

本節の内容は Wakabayashi \cite[\S2--3]{Wak} による。
本稿の視点からの再解釈を加えて述べる。
以下で引用する主張のうち「レベル $m$ $\leftrightarrow$ $p^{m+1}$-曲率」
（\cite[Def.~2.1.3]{Wak}）、「dormant なレベル $m$ $\iff$ $m+1$ 回 Frobenius 降下」
（\cite[Cor.~2.1.6]{Wak}）、「exponent は $\bbZ/p^{m+1}\bbZ$ に値をとる」
（\cite[Def.~3.3.1]{Wak}）は \texttt{handoff/app-session-2026-09-03.md} \S4 で
一次資料と照合されている。ただしこの照合は ar5iv 経由の二次確認であり
（同 \S0.3）、原論文の直接確認は付録 \ref{app:litcheck} の項目として残る。

\subsection{非対数の場合との対比}

\begin{proposition}[{\cite[Cor.~2.1.6]{Wak}}]\label{prop:nonlog}
$S\in\{R,K\}$ とする。関手 $\Xi^{\downarrow(m+1)}$, $\Xi^{\uparrow(m+1)}$ は圏同値
\[
  \bigl\{\text{dormant }\cD^{(m)}_S\text{-加群}\bigr\}
  \ \xrightarrow{\ \sim\ }\
  \bigl\{S^{(m+1)}\text{-加群}\bigr\}
\]
を与える。
\end{proposition}

すなわち非対数的な場合、dormant 対象は $F^{(m+1)*}$ の像に尽きる。
Wakabayashi は \cite[\S2.1 の Prop.~2.1.5 直後]{Wak} で
「しかし対数的な場合、すなわち $\breve\cD_R$-加群についてはこれは成り立たない」
と明示している。差を測る不変量が次の exponent である。

\subsection{階数 $1$ の分類}

\begin{proposition}[{\cite[Prop.~3.2.1]{Wak}}]\label{prop:rank1}
$a\in\bbZ/p^{m+1}\bbZ$ に対し $\tilde a$ を $0\le\tilde a<p^{m+1}$ なる代表とする。
\begin{equation}\label{eq:nabla_a}
  \nabla_{a,\langle j\rangle}(t^n):=q_j!\cdot\binom{n-\tilde a}{j}\cdot t^n
\end{equation}
は $R$ 上に $\Dbr{m}$-加群構造を定め、$p^{m+1}$-曲率は消える。
逆に、$E=R$ なる dormant $\Dbr{m}$-加群は
ただ一つの $a\in\bbZ/p^{m+1}\bbZ$ に対する $(R,\nabla_a)$ に同型である。
さらに
\[
  (R,\nabla_a)\otimes(R,\nabla_b)\cong(R,\nabla_{a+b}),\qquad
  (R,\nabla_a)^\vee\cong(R,\nabla_{-a}),
\]
\[
  \Hom\bigl((R,\nabla_a),(R,\nabla_b)\bigr)=
  \begin{cases}
    \{\mu_s\mid s\in R^{(m+1)}\} & (a=b),\\
    0 & (a\neq b).
  \end{cases}
\]
\end{proposition}

$m=0$ のとき $a\in\bbZ/p=\bbF_p$ であり、
$(R,\nabla_a)$ は \S\ref{ss:starting} の $\nabla=d+a'\dlog t$（$a'\in\bbF_p$）に対応する。
補題 \ref{lem:toral} が与える「留数は $\bbF_p$ 値」がここで
「exponent は $\bbZ/p$ 値」として再現される。

\begin{definition}[exponent]\label{def:exponent}
$(E,\nabla)$ を dormant $\Dbr{m}$-加群で $E$ が階数 $n$ の自由 $R$ 加群であるものとする。
$t$ 進完備化は
\[
  \xi:(\widehat E,\widehat\nabla)\xrightarrow{\ \sim\ }
  \bigoplus_{i=1}^n\bigl(k[\![t]\!],\widehat\nabla_{d_i}\bigr),
  \qquad d_i\in\bbZ/p^{m+1}\bbZ
\]
と分解する。多重集合 $[d_1,\dots,d_n]$ を $(E,\nabla)$ の \textbf{exponent} とよぶ。
\end{definition}

\subsection{$p$ 進展開としてのレベル}

次が本稿の描像の核心である。

\begin{proposition}[{\cite[Rem.~3.3.3]{Wak}}]\label{prop:padic}
$(E,\nabla)$ を dormant $\Dbr{m}$-加群、その exponent を $[d_1,\dots,d_n]$ とし、
$\tilde d_i=\sum_{l=0}^{m}p^l\,\tilde d_{il}$（$0\le\tilde d_{il}<p$）を $p$ 進展開とする。
$E^l:=\bigcap_{j=0}^{l-1}\operatorname{Ker}(\nabla_{\langle p^j\rangle})$、
$\overline\nabla^l:=\nabla_{\langle p^l\rangle}|_{E^l}$ とおくと、
$\breve\cD^{(0)}_{R^{(l)}}$-加群 $(E^l,\overline\nabla^l)$ のモノドロミーは対角化可能で、
その固有値は
\[
  -\tilde d_{1l},\ \dots,\ -\tilde d_{nl} \pmod p
\]
である。
\end{proposition}

\begin{corollary}[レベル上昇 $=$ 桁の追加]\label{cor:lifting-local}
局所的には、レベル $N$ の dormant 対象をレベル $N+1$ に持ち上げることは、
exponent $d\in\bbZ/p^{N}\bbZ$ の
$\bbZ/p^{N+1}\bbZ$ への持ち上げを選ぶことと同値であり、
各分岐点で $p$ 通りの選択肢がある。特に\textbf{局所的な障害は存在しない}。
\end{corollary}

\begin{remark}\label{rem:res-zero}
Wakabayashi \cite[Prop.~3.3.4]{Wak} により、次は同値である：
(a) 留数 $\Res(\nabla):=\Coker(\tau^{\downarrow\uparrow(m+1)}_{(E,\nabla)})$ が消える；
(b) exponent が $[0,\dots,0]$；
(c) $(E,\nabla)$ が非対数的な $\cD^{(m)}_R$-加群から $\eta_R$ を通じて来る。
すなわち\textbf{非対数的な理論とは exponent $=0$ の部分に他ならない}。
系 \ref{cor:lifting-local} の $p^r$ 通りの自由度は、
非対数的な世界では丸ごと消える。これが log と非 log の差の全てである。
\end{remark}

\subsection{cyclic vector と正則性}

\begin{proposition}[{\cite[Prop.~3.4.1]{Wak}}]\label{prop:cyclic}
dormant $\Dbr{m}$-加群 $(E,\nabla)$（$E$ は階数 $n$ の自由加群）に対し、
$m$-cyclic vector が存在することと、
exponent $d_1,\dots,d_n\in\bbZ/p^{m+1}\bbZ$ が互いに相異なることは同値である。
\end{proposition}

$m$-cyclic vector をもつ $\cD^{(m)}$-加群は局所的に $\mathrm{GL}_n$-oper に他ならない
\cite[Prop.~4.2.3]{Wak}。
したがって命題 \ref{prop:cyclic} は
「oper $\Rightarrow$ 留数は正則」という古典的命題の高次レベル版である。
$n=2$ の場合、$d_1\neq d_2$ は次節の $\rho\neq0$ に対応する。

%%%%%%%%%%%%%%%%%%%%%%%%%%%%%%%%%%%%%%%%%%%%%%%%%%%%%%%%%%%%%%%%%%%%%%%%%%%%%%%
\section{格子の取り替えと次数障害}\label{sec:lattice}

\subsection{Hecke 変形は局所的には無力}

\begin{lemma}\label{lem:hecke}
$\nabla_\theta(x^{-n}f)=x^{-n}(\theta f+(a-n)f)$ より、
$D$ に沿った格子の取り替えは exponent を $a\mapsto a-n$（$n\in\bbZ$）だけずらす。
特にレベル $N$ では exponent は $\bbZ/p^N$ 全体を動きうる。
\end{lemma}

したがって階数 $1$ では、局所的に exponent を常に $0$ にできる。
格子に依存しない不変量は $\nabla|_{X\setminus D}$ の同型類のみである。
「持ち上がらない」ことを主張するには、
\textbf{格子が剛的である}状況、すなわち oper のような構造が必要になる。

\subsection{階数 $1$・曲線の場合の完全な答え}

\begin{proposition}\label{prop:rank1curve}
$X$ を固有滑らかな曲線、$D=\sum_{i=1}^r p_i$（$r\ge1$）、
$(\cL,\nabla)$ を階数 $1$ の dormant 対数接続とする。
このとき $(\cL,\nabla)$ は任意のレベルに持ち上がる。
すなわち stratification をもつ。
\end{proposition}

\begin{proof}
命題 \ref{prop:rank1} と系 \ref{cor:lifting-local} により、
レベル $N$ をもつことは、ある直線束 $\cN$ と整数 $j_i$（$0\le j_i<p^N$、
$j_i$ は与えられた exponent の持ち上げ）に対し
\[
  \cL\cong F^{(N)*}\cN\otimes\cO\Bigl(\sum_i j_i\,p_i\Bigr)
\]
と書けることと同値。次数をとると
$\deg\cL=p^N\deg\cN+\sum_i j_i$、すなわち
\[
  \sum_i j_i\equiv\deg\cL \pmod{p^N} .
\]
$r\ge1$ なので $j_1$ を $\bmod\,p^N$ で自由に選べ、この合同式は解ける。
残りは $\Pic^0(X)$ が $p$ 可除であること（$X$ は固有滑らかな曲線）から従う。
\end{proof}

\begin{remark}[主張の地位]\label{rem:rank1curve-status}
命題 \ref{prop:rank1curve} は \emph{[本研究・部分証明]} 扱いとする。
局所分類（命題 \ref{prop:rank1} $=$ Wakabayashi \cite[Prop.~3.2.1]{Wak}）の大域化、
すなわち「各分岐点で独立に選んだ exponent の持ち上げが大域的な直線束 $\cN$ と
両立する」ことの細部を精査していない（\texttt{handoff/level-chat.md} \S3.2）。
\end{remark}

\begin{remark}
$r=0$ のときは $j_i$ の自由度が無く、
$\deg\cL\equiv0\bmod p^N$ が必要になる。
これが古典的な障害であり、次の例を与える。
\end{remark}

\begin{example}[ちょうどレベル $N$ の例、非対数]\label{ex:exactN}
$X$ を種数 $\ge1$ の曲線、$D=\emptyset$、
$\cM$ を $p\nmid\deg\cM$ なる直線束とし、$\cL:=F^{(N)*}\cM$（標準接続つき）とする。
$\cL$ は dormant レベル $N$ であるが、
$\deg\cL=p^N\deg\cM$ は $p^{N+1}$ で割れないのでレベル $N+1$ ではない。
すなわち\textbf{ちょうどレベル $N$}である。
\end{example}

\subsection{一般の次数障害}

\begin{proposition}[次数障害]\label{prop:degree}
$(X,D)$ を $r$ 点付き固有滑らかな曲線（種数 $g$、$D=\sum_{i=1}^r\sigma_i$）とし、
$(\cF,\nabla,\cL)$ を dormant レベル $N$ の $\mathrm{GL}_2$-oper とする。
$\det\cF$ の exponent を $e_i\in\bbZ/p^N\bbZ$、
$\tilde e_i\in[0,p^N)$ をその代表とすると
\[
  \det\cF\cong F^{(N)*}\cM\otimes\cO\Bigl(\sum_i\tilde e_i\,\sigma_i\Bigr)
\]
なる直線束 $\cM$ が存在し、
\begin{equation}\label{eq:degcong}
  \boxed{\ \sum_{i=1}^r\tilde e_i\equiv\deg\det\cF \pmod{p^N}\ }
\end{equation}
が成り立つ。さらに Kodaira--Spencer 同型により
$\cL^{\otimes2}\cong\Omega^1_X(\log D)\otimes\det\cF$、
したがって
\[
  2\deg\cL=(2g-2+r)+\deg\det\cF .
\]
\end{proposition}

\begin{proof}
$\det\cF$ は階数 $1$ の dormant $\Dlog{N-1}$-加群なので
命題 \ref{prop:rank1} と remark \ref{rem:res-zero} により
$\det\cF\otimes\cO(-\sum\tilde e_i\sigma_i)$ の留数が消え、
命題 \ref{prop:nonlog}（$=$ \cite[Cor.~2.1.6]{Wak}）の\emph{大域版}%
\NeedsCheck{dormant $\cD^{(m)}$-加群に対する Frobenius 降下 (Cor 2.1.6) の
大域的な形の正確な参照が未確定}%
により $F^{(N)*}$ の像に入る。
次数を取れば \eqref{eq:degcong} を得る。
KS 同型は oper の定義から。
\end{proof}

\begin{corollary}[持ち上げの障害の所在]\label{cor:obstruction}
レベル $N$ から $N+1$ への持ち上げは、
各 $e_i$ を $\bbZ/p^{N+1}$ に持ち上げて \eqref{eq:degcong} を
$\bmod\,p^{N+1}$ で満たすことに帰着する。
\begin{itemize}[leftmargin=1.6em]
\item $r\ge1$：$\tilde e_1$ の持ち上げだけで $\bmod\,p$ の自由度があるので、
      合同式は常に解ける。
\item $r=0$：自由度が無く、$\deg\det\cF\equiv0\bmod p^{N}$ が必要。
      これが $\Pic$ の $p$ 可除性に関する障害を与える。
\end{itemize}
\end{corollary}

\begin{remark}
ただし \eqref{eq:degcong} は必要条件であって十分条件ではない。
oper の存在にはさらに幾何的な条件（\S\ref{sec:opers} の三角不等式に相当するもの）
が課される。$\sP=(\bbP^1;0,1,\infty)$ の場合には
定理 \ref{thm:M03}（仮定 (W10.3) のもと）でこれを扱う。
一般の $(g,r)$ で十分性を論じることは未解決である（\S\ref{sec:open}、問題 \ref{q:degree-suff}）。
\end{remark}

%%%%%%%%%%%%%%%%%%%%%%%%%%%%%%%%%%%%%%%%%%%%%%%%%%%%%%%%%%%%%%%%%%%%%%%%%%%%%%%
\section{$\mathrm{PGL}_2$-oper、radius、被覆}\label{sec:opers}

\subsection{定義と radius}

$\sX=(X,\{\sigma_i\}_{i=1}^r)$ を $r$ 点付き滑らかな曲線とする。

\begin{definition}[{\cite[Def.~4.2.1, 4.2.5]{Wak}}]
レベル $N$ の $\mathrm{GL}_2$-oper とは、
$\Dlog{N-1}$-加群 $(\cF,\nabla)$（$\cF$ は階数 $2$）と直線部分束 $\cL\subset\cF$ の組で、
合成
\[
  \cL\hookrightarrow\cF\xrightarrow{\ \overline\nabla\ }
  \Omega_{X^{\log}}\otimes\cF\twoheadrightarrow\Omega_{X^{\log}}\otimes(\cF/\cL)
\]
が同型となるものをいう。
$\mathrm{PGL}_2$-oper は、階数 $1$ の $\Dlog{N-1}$-加群による捻りを法とした同値類である。
$p^N$-曲率が消えるとき dormant という。
${}^N\Op^{\Zzz}_{2,\sX}$ で dormant $\mathrm{PGL}_2$-oper（レベル $N$）の集合を表す。
\end{definition}

\begin{definition}[radius, {\cite[Def.~4.3.2, Rem.~4.3.1]{Wak}}]
$\sigma_i$ における exponent $[d_1,d_2]\subset\bbZ/p^N\bbZ$ の
$\Delta$（対角）による類を radius とよぶ。$n=2$ のとき
$a\mapsto[a,-a]$ により
\[
  \fS_2\backslash(\bbZ/p^N\bbZ)^{\times2}/\Delta\ \cong\ (\bbZ/p^N\bbZ)/\{\pm1\}
\]
と同一視し、$\rho_i=\tfrac12(d_1-d_2)\in(\bbZ/p^N\bbZ)/\{\pm1\}$ と書く。
\end{definition}

\begin{proposition}[{\cite[Rem.~4.3.3]{Wak}}, {\cite[Chap.~II, Prop.~1.4, 1.5]{Moc}}]
\label{prop:radius-unit}
$N=1$ のとき radius は $\bbF_p$ に属し、
${}^1\Op^{\Zzz}_{2,\sX,\vec\rho}$ は
$\rho_i\in\bbF_p^\times/\{\pm1\}$（すべての $i$）でない限り空である。
\end{proposition}

これは補題 \ref{lem:toral} と命題 \ref{prop:cyclic} の組み合わせでもある：
$\psi=0$ より留数は半単純、oper より正則、
ゆえに二つの指数は相異なり $\rho\in\bbF_p^\times$。
この別証明は $N=1$ について本稿内で自足しており、Mochizuki の引用に依存しない。
\NeedsCheck{Mochizuki \cite[Chap.~II, Prop.~1.4, 1.5]{Moc} の正確な主張は
未確認（Wakabayashi 経由の引用のみ。\texttt{handoff/level-chat.md} \S7 \#8）。}

\subsection{被覆との対応}

\begin{theorem}[{\cite[Thm.~C]{Wak}}, $N=1$ は {\cite[Introduction, Thm.~1.3]{Moc}}]
\label{thm:covering}
$\sP=(\bbP^1;[0],[1],[\infty])$ とする。
$\mathrm{Cov}^{\mathrm{tame}}_{\sP}$ を、
分岐点の集合が $\{[0],[1],[\infty]\}$ に一致し、
分岐指数 $\lambda_0,\lambda_1,\lambda_\infty$ がすべて奇数で
$\lambda_0+\lambda_1+\lambda_\infty<2p^N$ を満たす
有限・分離的・従順分岐被覆 $\phi:\bbP^1\to\bbP^1$ の
（$\mathrm{PGL}_2(k)$ を左から合成する同値による）類の集合とする。
このとき標準的な全単射
\[
  \Upsilon_{\sP}:\mathrm{Cov}^{\mathrm{tame}}_{\sP}
  \xrightarrow{\ \sim\ }{}^N\Op^{\Zzz}_{2,\sP}
\]
があり、radius は $\bigl(\tfrac12\overline{\lambda_0},
\tfrac12\overline{\lambda_1},\tfrac12\overline{\lambda_\infty}\bigr)$ の像に一致する。
\end{theorem}

\subsection{被覆側の初等的な言い換え}

$p$-曲率が消えることは、$\cE^\nabla$ が $\cO_{X'}$ 上階数 $2$ であること、
すなわち対応する $2$ 階線型微分方程式が
$k(x)$ 内に $k(x^p)$ 上一次独立な解 $y_1=f$, $y_2=g$ をもつことと同値である。
このとき $\phi=f/g$ が定理 \ref{thm:covering} の被覆であり、
Wronskian $W=f'g-fg'$ が $\phi'=W/g^2$ を通じて分岐を支配する：
\begin{equation}\label{eq:wronskian}
  \operatorname{ord}_{x_0}W=\lambda_{x_0}-1,\qquad
  \deg W=2d-2-(\lambda_\infty-1),\qquad d=\deg\phi .
\end{equation}
Riemann--Hurwitz から $\sum_x\lambda_x=2d+1$、したがって
$\lambda_{x}\le d$ は
\begin{equation}\label{eq:triangle}
  \lambda_0<\lambda_1+\lambda_\infty
  \quad\text{（および巡回置換）}
\end{equation}
と同値である。すなわち\textbf{厳密三角不等式は被覆の次数条件に他ならない}。

\subsection{数え上げ}

\begin{lemma}\label{lem:lambda-unique}
$p^N$ は奇数なので、$0<a<p^N$ なる $a$ に対し $\{a,p^N-a\}$ のちょうど一方が奇数である。
したがって radius $\rho\in(\bbZ/p^N\bbZ)^\times/\{\pm1\}$ に対し、
$\lambda\equiv\pm2\rho\pmod{p^N}$、$\lambda$ 奇数、$0<\lambda<p^N$
なる $\lambda$ は一意に定まる。
\end{lemma}

\begin{proposition}\label{prop:count}
レベル $N=1$ のとき
\begin{equation}\label{eq:count1}
  \#\,{}^1\Op^{\Zzz}_{2,\sP}
  =\#\Bigl\{(\lambda_0,\lambda_1,\lambda_\infty)\ \Big|\
  \begin{array}{l}
  \lambda_x\ \text{奇数},\ 0<\lambda_x<p,\\[2pt]
  \text{\eqref{eq:triangle} の三角不等式},\ \textstyle\sum_x\lambda_x<2p
  \end{array}\Bigr\}
  =\frac{p(p^2-1)}{24}
  =\frac{|\mathrm{PSL}_2(\bbF_p)|}{12} .
\end{equation}
一般のレベル $N$ については、命題 \ref{prop:Cdagger} により
$\#\,{}^N\Op^{\Zzz}_{2,\sP}=\#\,{}^{\dagger}C_N$ である。
$N\ge2$ では、素朴な条件
（$\lambda_x$ 奇、$p\nmid\lambda_x$、三角、$\sum\lambda_x<2p^N$）のみでは
正しい数え上げにならない（注意 \ref{rem:erratum}）。
\end{proposition}

\begin{proof}[証明（前半、および後半の還元）]
前半（左辺 $=$ $\lambda$ の集合）は定理 \ref{thm:covering} と補題 \ref{lem:lambda-unique}
による。$N=1$ では条件 \ref{def:Cdagger}(b) が空なので補題 \ref{lem:dict} により
この集合は ${}^{\dagger}C_1$ と一致する。
後半（$\#\,{}^{\dagger}C_1=p(p^2-1)/24$）は $p\in\{3,5,7,11,13\}$ で計算機により
確認済み（表 \ref{tab:count}、付録 \ref{app:num}）だが、一般 $p$ に対する
証明は与えていない。
\GAP{$\#\,{}^{\dagger}C_1=n(n+1)(2n+1)/6$（$n=(p-1)/2$、四角錐数；なお
$n(n+1)(2n+1)/6=p(p^2-1)/24$ は $n=(p-1)/2$ の代入により恒等的に成立する）
の組合せ論的証明は未記述。純粋に組合せ論の問題
（\texttt{handoff/level-chat.md} \S6.3）。}
\end{proof}

\subsection{Wakabayashi の集合 ${}^{\dagger}C_N$}\label{ss:Cdagger}

$N=1$ の場合の記述は、次の形で任意のレベルに一般化される。
$a\in\bbZ_{\ge0}$ と $N'\in\bbZ_{>0}$ に対し $[a]_{N'}:=a\bmod p^{N'}$ と書く。

\begin{definition}[{\cite[(10.6)]{Wak2}}]\label{def:Cdagger}
${}^{\dagger}C_N$ を、次の二条件を満たす非負整数の三つ組 $(s_1,s_2,s_3)$ 全体とする。
\begin{enumerate}[label=(\alph*),leftmargin=2em]
\item $\displaystyle\sum_{i=1}^3 s_i\le p^N-2$ かつ $|s_2-s_3|\le s_1\le s_2+s_3$;
\item 任意の $0<N'<N$ に対し、$s_i'\in\{[s_i]_{N'},\ p^{N'}-1-[s_i]_{N'}\}$ $(i=1,2,3)$ を
      適当に選んで $\sum_i s_i'\le p^{N'}-2$ かつ $|s_2'-s_3'|\le s_1'\le s_2'+s_3'$
      とできる。
\end{enumerate}
\end{definition}

\begin{proposition}[{\cite[Prop.~10.1.4, 10.3.3]{Wak2}}]\label{prop:Cdagger}
dormant $\mathrm{PGL}_2$-oper（レベル $N$）on $\sP$ の集合は ${}^{\dagger}C_N$ で
径数付けられる。すなわち $\#\,{}^N\Op^{\Zzz}_{2,\sP}=\#\,{}^{\dagger}C_N$。
\end{proposition}

\NeedsCheck{命題 \ref{prop:Cdagger} は Wakabayashi \cite{Wak2} \S10 の結果である。
本稿では \cite{Wak2} \S10.1--10.3（特に Prop.~10.3.3 の径数付け）を
原論文で確認していない（\texttt{handoff/level-chat.md} \S7 \#2、
\texttt{handoff/app-session-2026-09-03.md} の L2）。また arXiv:2209.08528 の
題名について書誌上の不一致がある（付録 \ref{app:litcheck}）。}

\begin{lemma}[辞書]\label{lem:dict}
$\lambda_i=2s_i+1$ とおくと、条件 \ref{def:Cdagger}(a) は
「$\lambda_i$ は奇数、$\sum\lambda_i\le 2p^N-1$、厳密三角不等式 \eqref{eq:triangle}」
と同値である。特に $s_i=0$ は $\lambda_i=1$（不分岐）に対応し、
定理 \ref{thm:covering} の $\mathrm{Cov}^{\mathrm{tame}}_{\sP}$ は
$\lambda_x=1$ を許す読み方が正しい。
\end{lemma}

\begin{proof}
$\sum\lambda_i=2\sum s_i+3\le 2(p^N-2)+3=2p^N-1<2p^N$。
また $\lambda_1<\lambda_2+\lambda_3\iff 2s_1+1<2s_2+2s_3+1\iff s_1\le s_2+s_3$
（$s_i$ は整数）。$|s_2-s_3|\le s_1$ は残り二つの三角不等式に同値。
\end{proof}

\begin{remark}[前稿 ver.\,0.3 の数値の誤りの訂正]\label{rem:erratum}
条件 \ref{def:Cdagger}(b) は $N\ge2$ で真に効く。これを落とした素朴な条件
（$\lambda_i$ 奇、$p\nmid\lambda_i$、三角、$\sum\lambda_i<2p^N$）による数え上げは
$N=1$ では正しいが（(b) が空条件）、$N\ge2$ では過大評価となる。
計算機による確認（付録 \ref{app:num}）では、$N=2$ で
\[
  \#\,{}^{\dagger}C_2\ (\text{正}):\quad 11,\ 225,\ 1666,\ 24805,\ 67431
  \qquad(p=3,5,7,11,13),
\]
\[
  \text{素朴な条件による値}:\quad 11,\ 345,\ 3122,\ 55605,\ 158431 .
\]

前稿 ver.\,0.3 の表では $N=2$ の値を $p=5$ で $610$、$p=7$ で $2262$ と
「計算機による直接列挙」として記載していた。しかしこの列挙は実際には
実行されておらず、値も誤りである（\texttt{handoff/level-chat.md} \S5.4, 付録B、
本稿付録 \ref{app:errata}）。正しい値は上記および表 \ref{tab:count} の通りで、
素朴な条件による過大評価値 $345,\ 3122$ とも一致しない。
\end{remark}

\begin{table}[htbp]
\centering
\caption{$\#\,{}^N\Op^{\Zzz}_{2,\sP}=\#\,{}^{\dagger}C_N$（定義 \ref{def:Cdagger} による直接列挙）。
$N\ge2$ の値は \texttt{computations/count\_dagger\_C.py}（および回収した
\texttt{computations/check.py}）を本稿の作業環境で実行して得たもの（付録 \ref{app:num}）。}
\label{tab:count}
\begin{tabular}{lrrrrr}
\toprule
$p$ & 3 & 5 & 7 & 11 & 13\\
\midrule
$N=1$ & 1 & 5 & 14 & 55 & 91\\
$p(p^2-1)/24$ & 1 & 5 & 14 & 55 & 91\\
$N=2$ & 11 & 225 & 1666 & 24805 & 67431\\
\bottomrule
\end{tabular}
\end{table}

さらに $p=3$ では $\#\,{}^{\dagger}C_3=121=11^2$ を確認した（付録 \ref{app:num}）。

\begin{remark}[$g=2$ との一致]\label{rem:genus2}
$(p^3-p)/24$ は、一般な種数 $2$ の曲線上の dormant $\mathrm{PGL}_2$-oper の個数
（Mochizuki, Lange--Pauly, Osserman）と一致する \cite[(1.2)]{Wak2}。
これは偶然ではないと考えられる：種数 $2$ の三価グラフはテータグラフ
（頂点 $2$、辺 $3$）であり、両頂点が同じ辺三つ組を見るので、
balanced edge numbering の集合は ${}^{\dagger}C_N$ そのものになる \cite[Thm.~E]{Wak2}。
すなわち $\deg\bigl(\Pi_{2,N,\emptyset,2,0}\bigr)=\#\,{}^{\dagger}C_N$。
\NeedsCheck{\cite[Thm.~E, (1.2)]{Wak2} は原論文で未確認（付録 \ref{app:litcheck}）。}
\end{remark}

%%%%%%%%%%%%%%%%%%%%%%%%%%%%%%%%%%%%%%%%%%%%%%%%%%%%%%%%%%%%%%%%%%%%%%%%%%%%%%%
\section{$p=5$ の明示計算}\label{sec:p5}

$p=5$、$\sP=(\bbP^1;0,1,\infty)$ とする。
$\bbF_5^\times/\{\pm1\}=\{[1],[2]\}$ であり、
$\lambda\mapsto\rho=\lambda/2=3\lambda$ により
$\lambda=1\mapsto[2]$、$\lambda=3\mapsto[1]$。
補題 \ref{lem:dict} により $\lambda_i=2s_i+1$、$(s_i)\in{}^{\dagger}C_1$、
すなわち $\sum s_i\le3$ かつ三角不等式。該当する $(s_1,s_2,s_3)$ は
\[
  (0,0,0),\quad(1,1,0),\quad(1,0,1),\quad(0,1,1),\quad(1,1,1)
\]
の $5$ 個であり、$\lambda$ に直すと
$(1,1,1),(3,3,1),(3,1,3),(1,3,3),(3,3,3)$。
除外されるのは $(1,1,3)$ 型（$3<1+1$ に反する）。
radius 三つ組 $8$ 通りのうち $(2,2,1)$, $(2,1,2)$, $(1,2,2)$ は実現されない。
この $5$ 個の $(s_i)$ は \texttt{computations/count\_dagger\_C.py} の出力
（$p=5$, $N=1$ の全リスト）と一致する（付録 \ref{app:num}）。

\begin{table}[htbp]
\centering
\caption{$p=5$ における dormant $\mathrm{PGL}_2$-oper の完全リスト（レベル $1$）。
Wronskian と各 ODE は \texttt{computations/verify\_p5.py}
（$=$ \texttt{handoff/verify\_p5.py}）で $\bbF_5$ 上検証済み（付録 \ref{app:num}）。}
\label{tab:p5}
\small
\begin{tabular}{clcclc}
\toprule
\# & $(\lambda_0,\lambda_1,\lambda_\infty)$ & $(\rho_0,\rho_1,\rho_\infty)$ & $d$
   & 被覆 $\phi$ / ODE & ${}_2F_1(a,b;c)$\\
\midrule
1 & $(1,1,1)$ & $(2,2,2)$ & 1 & $\phi=x$;\quad $y''=0$ & $(4,0;0)$\\[2pt]
2 & $(3,1,3)$ & $(1,2,1)$ & 3 & $\phi=x^3$;\quad $xy''-2y'=0$ & ---\\[2pt]
3 & $(1,3,3)$ & $(2,1,1)$ & 3 & $\phi=(x-1)^3$;\quad $(x-1)y''-2y'=0$ & ---\\[2pt]
4 & $(3,3,1)$ & $(1,1,2)$ & 3 & $\phi=\dfrac{x^3}{(x-1)^3}$;\quad
      $x(x-1)y''+(x+2)y'+y=0$ & $(2,3;3)$\\[8pt]
5 & $(3,3,3)$ & $(1,1,1)$ & 4 & $\phi=\dfrac{x^3(2x+1)}{x+2}$;\quad
      $x(x-1)y''+(x+2)y'+4y=0$ & $(1,4;3)$\\
\bottomrule
\end{tabular}
\end{table}

\subsection{\#5 の構成（詳細）}\label{ss:no5}

$\lambda=(3,3,3)$、$d=4$。\eqref{eq:wronskian} より $W=c\,x^2(x-1)^2$（次数 $4$）で、
$2d-2=6$ からの次数落ち $2$ が $\lambda_\infty-1=2$ に対応する。
$\infty$ での指数が $3$ なので $\deg f=4$, $\deg g=1$。
$g=x-t$、$f=\sum_{i=0}^4a_ix^i$ とおくと
\[
  W=f'g-fg'=\sum_i(i-1)a_ix^i-t\sum_i ia_ix^{i-1} .
\]
係数比較（$\bbF_5$）：
\[
  \begin{array}{ll}
  x^4:&3a_4=c\\
  x^3:&2a_3-4ta_4=-2c/3\\
  x^2:&a_2-3ta_3=c/3\\
  x^1:&-2ta_2=0\\
  x^0:&-a_0-ta_1=0 .
  \end{array}
\]
$a_4=2$（よって $c=6=1$）と正規化する。
$x^1$ の式から $t=0$ または $a_2=0$。$t=0$ は $g$ の零点が分岐点 $x=0$ と衝突するので
$a_2=0$ を採る。$x^3$ の式から $a_3=4-t$、
$x^2$ の式 $-3ta_3=1$ に代入して
\[
  t(4-t)=3\ \Longleftrightarrow\ t^2-4t+3=(t-1)(t-3)=0 .
\]
$t=1$ は分岐点 $x=1$ と衝突するので $t=3$、$a_3=1$。
$a_1$ は $f\mapsto f+a_1g$（$=\phi\mapsto\phi+a_1$、Möbius）の自由度なので $a_1=a_0=0$。
よって
\[
  f=2x^4+x^3,\qquad g=x+2 .
\]
検算：$f'=3x^3+3x^2$、
\[
  W=(3x^3+3x^2)(x+2)-(2x^4+x^3)=x^4+3x^3+x^2=x^2(x-1)^2 .\quad\checkmark
\]

微分方程式は
$\det\bigl(\begin{smallmatrix}y&f&g\\ y'&f'&g'\\ y''&f''&g''\end{smallmatrix}\bigr)=0$、
すなわち
\[
  (fg'-f'g)\,y''+(f''g-fg'')\,y'+(f'g''-f''g')\,y=0 .
\]
$f''=4x^2+x$、$g''=0$ より係数はそれぞれ
$4x^2(x-1)^2$、$4x(x-1)(x+2)$、$x(x-1)$。
$x(x-1)$ で割り $4$ 倍して
\[
  x(x-1)y''+(x+2)y'+4y=0 .
\]
検算：$y=x+2$ で $(x+2)+4(x+2)=5(x+2)=0$；
$y=2x^4+x^3$ で $5x^4+5x^3=0$。\quad$\checkmark$

超幾何形 $x(1-x)y''+[c-(a+b+1)x]y'-aby=0$ と比較すると
$c=3$、$a+b=0$、$ab=4$、すなわち $(a,b,c)=(1,4;3)$。
指数は $0$ で $\{0,1-c\}=\{0,3\}$、$1$ で $\{0,c-a-b\}=\{0,3\}$、
$\infty$ で $\{a,b\}=\{1,4\}$（差 $3$）。
三点とも指数差 $3$、すなわち $\lambda=(3,3,3)$。\quad$\checkmark$

\subsection{\#4 の構成}

$\phi=x^3/(x-1)^3$、$f=x^3$、$g=(x-1)^3$。
$W=3x^2(x-1)^3-3x^3(x-1)^2=-3x^2(x-1)^2=2x^2(x-1)^2$、
$\deg W=4=2d-2$ より $\infty$ では不分岐（$\lambda_\infty=1$）。
同様の計算で
\[
  x(x-1)y''+(x+2)y'+y=0,\qquad (a,b,c)=(2,3;3),
\]
指数は $0,1$ で差 $3$、$\infty$ で $\{2,3\}$（差 $1$）。
すなわち $\lambda=(3,3,1)$。

\subsection{レベル $2$ への持ち上げ}

補題 \ref{lem:nesting} により、上の $5$ 個の $(s_i)$ はそのまま
${}^{\dagger}C_2$ の元でもある（$\sum s_i\le3\le p^2-2=23$、
条件 \ref{def:Cdagger}(b) は $N'=1$ で $s_i'=s_i$ と選べる。付録 \ref{app:num} で確認）。
したがって命題 \ref{prop:Cdagger} と仮定 (W10.3)（\S\ref{sec:M03}）のもとで
$5$ 個すべてがレベル $2$ に持ち上がる。レベル $2$ の総数は $\#\,{}^{\dagger}C_2=225$
（表 \ref{tab:count}）。

%%%%%%%%%%%%%%%%%%%%%%%%%%%%%%%%%%%%%%%%%%%%%%%%%%%%%%%%%%%%%%%%%%%%%%%%%%%%%%%
\section{持ち上げの障害は行列式に尽きる}\label{sec:M03}

本節の主結果（定理 \ref{thm:M03}、系 \ref{cor:general}）は、次の仮定のもとでの
\emph{条件付き定理} である。

\begin{quote}
\textbf{仮定 (W10.3).}\ 命題 \ref{prop:Cdagger} の径数付け
$\#\,{}^N\Op^{\Zzz}_{2,\sP}=\#\,{}^{\dagger}C_N$ のもとで、レベル降下
$\mathrm{res}_N:{}^{N+1}\Op^{\Zzz}_{2,\sP}\to{}^{N}\Op^{\Zzz}_{2,\sP}$ は
${}^{\dagger}C_N$ の言葉で
$s\mapsto s'\in\{[s]_N,\ p^N-1-[s]_N\}$（各成分ごとの $\pm$ の選択）として
実現される。
\end{quote}

\NeedsCheck{仮定 (W10.3) は Wakabayashi \cite{Wak2} \S10.3
（特に Prop.~10.3.3）で確認すべき事項であり、本稿では原論文を確認していない
（\texttt{handoff/level-chat.md} \S3.3, \S7 \#2、
\texttt{handoff/app-session-2026-09-03.md} の L2）。
条件 \ref{def:Cdagger}(b) の形からほぼ自明に見えるが、径数付けそのものとの
両立は別個の主張である。}

\subsection{${}^{\dagger}C_N$ の入れ子性}

\begin{lemma}\label{lem:nesting}
${}^{\dagger}C_N\subseteq{}^{\dagger}C_{N+1}$。
\end{lemma}

\begin{proof}
$(s_i)\in{}^{\dagger}C_N$ とする。$\sum s_i\le p^N-2\le p^{N+1}-2$ かつ三角不等式より
条件 \ref{def:Cdagger}(a) は $N+1$ でも成立する。
条件 (b) について、$N'=N$ のときは $s_i<p^N$ ゆえ $[s_i]_N=s_i$ であり、
$s_i'=s_i$ と選べば要求は $(s_i)$ に対する条件 (a)（レベル $N$）そのものである。
$N'<N$ のときは $(s_i)\in{}^{\dagger}C_N$ の条件 (b) をそのまま使えばよい。
\end{proof}

\begin{remark}
補題 \ref{lem:nesting} は $p\in\{3,5\}$、$N\in\{1,2\}$ について計算機でも確認した
（付録 \ref{app:num}：${}^{\dagger}C_N\setminus{}^{\dagger}C_{N+1}$ は空）。
\end{remark}

\subsection{$\sP$ 上での全射性}

\begin{theorem}[仮定 (W10.3) のもとで]\label{thm:M03}
仮定 (W10.3) を仮定する。$\sP=(\bbP^1;0,1,\infty)$ に対し、制限写像
\[
  \mathrm{res}_N:\ {}^{N+1}\Op^{\Zzz}_{2,\sP}\longrightarrow{}^{N}\Op^{\Zzz}_{2,\sP}
\]
は全射である。特に $\sP$ 上のすべての dormant $\mathrm{PGL}_2$-oper は
任意に高いレベルに持ち上がり、stratification に延びる。
\end{theorem}

\begin{proof}
命題 \ref{prop:Cdagger} により両辺は ${}^{\dagger}C_{N+1}$、${}^{\dagger}C_N$ と
同一視される。仮定 (W10.3) により $\mathrm{res}_N$ は
$s\mapsto s'\in\{[s]_N,\ p^N-1-[s]_N\}$ で与えられる。
補題 \ref{lem:nesting} の包含 ${}^{\dagger}C_N\subseteq{}^{\dagger}C_{N+1}$ の各元は
$[s]_N=s$ を満たすので、この包含が $\mathrm{res}_N$ の切断を与える。
\end{proof}

\begin{remark}[主張の地位]\label{rem:M03-status}
定理 \ref{thm:M03} は仮定 (W10.3) のもとで完全な証明をもつ
\emph{[本研究・条件付き定理]}。仮定 (W10.3)（および参照した \cite{Wak2}
の径数付け）が原論文で確認されれば無条件の定理に格上げされる。
これが本研究の主要主張（$\sP$ 上に反例が存在しないこと）が確定するか否かの
分岐点である（\texttt{handoff/app-session-2026-09-03.md} \S11-2）。
\end{remark}

\subsection{一般の $(g,r)$ へ}

\begin{corollary}[仮定 (W10.3) と {\cite[Thm.~E]{Wak2}} のもとで]\label{cor:general}
仮定 (W10.3) を仮定する。$\mathbb{G}$ を型 $(g,r)$ の三価 clutching data、
$X_{\mathbb{G}}$ を対応する totally degenerate curve とする。
balanced $(p,N)$-edge numbering の集合について
$\mathrm{Ed}_{p,N,\mathbb{G}}\subseteq\mathrm{Ed}_{p,N+1,\mathbb{G}}$ が成り立ち、
$X_{\mathbb{G}}$ 上の制限写像は全射である。
\end{corollary}

\begin{proof}
$\mathrm{Ed}$ の各頂点条件は「接する三辺の番号が ${}^{\dagger}C_N$ に属する」であり
\cite[Thm.~E]{Wak2}、補題 \ref{lem:nesting} を各頂点に適用すればよい。
全射性は定理 \ref{thm:M03} と同じ議論を各頂点で行う。
\end{proof}

\NeedsCheck{系 \ref{cor:general} の証明は、totally degenerate curve 上の
dormant oper が頂点ごとの ${}^{\dagger}C_N$-データで径数付けられること
（\cite[Thm.~E]{Wak2}）を用いる。\cite{Wak2} \S9--10 は本稿では未確認
（付録 \ref{app:litcheck}）。}

\begin{remark}[主張の地位]\label{rem:no-pgl2-cex}
系 \ref{cor:general} は、当初の主要予想（\S\ref{sec:open} の問題 \ref{q:main}）に対する
\textbf{否定的な}答えを、仮定 (W10.3) のもとで、かつ totally degenerate な曲線に限って
与える：$\mathrm{PGL}_2$-oper については、$r=0$ を含めどの $(g,r)$ でも
「ちょうどレベル $N$」の反例は存在しない。
\emph{[本研究・条件付き系]}。
\Todo{一般の smooth な曲線上で $\mathrm{res}_N$ が全射であることを、
generic étaleness \cite[Thm.~C]{Wak2} から直接導く議論を書く。
現状は totally degenerate な曲線と、それに基づく次数の単調性
（$\deg\Pi_{2,N,\rho,g,r}\le\deg\Pi_{2,N+1,\rho',g,r}$、$\rho'$ は $\rho$ の
自明な持ち上げ）までしか押さえていない。}
\end{remark}

\subsection{障害は行列式に局在する}

$\mathrm{PGL}_2$ で（上記の条件のもと）障害が消える一方、$\mathrm{GL}_1$
（例 \ref{ex:exactN}）では「ちょうどレベル $N$」が実現する。この差は次のように
整理できる。

\begin{theorem}[障害の局在]\label{thm:det}
$(\cF,\nabla,\{\cF^j\})$ を dormant レベル $N$ の $\mathrm{GL}_n$-oper とし、
$\mathscr{F}^{\spadesuit}$ をその $\mathrm{PGL}_n$-類とする。
$\mathscr{F}^{\spadesuit}$ がレベル $N+1$ に持ち上がり、かつ $p\nmid n$ と仮定する。
このとき $(\cF,\nabla,\{\cF^j\})$ 自身がレベル $N+1$ に持ち上がるための必要十分条件は、
階数 $1$ の dormant $\Dlog{N-1}$-加群 $(\det\cF,\det\nabla)$ が
レベル $N+1$ に持ち上がることである。
\end{theorem}

\begin{proof}[証明の概略]
$\mathrm{GL}_n$-oper は $\mathrm{PGL}_n$-oper と、階数 $1$ の dormant $\cD^{(N-1)}$-加群
による捻りの差しかない（\S\ref{sec:opers} の定義）。
持ち上げ $\widetilde{\mathscr{F}^{\spadesuit}}$ を一つ選び、その代表 $\mathrm{GL}_n$-oper を
$\widetilde{\cF}$ とすると、$\det\widetilde{\cF}$ と $\det\cF$ は
レベル $N$ で同型な階数 $1$ の対象を与える。
必要な捻り $(\cN,\nabla_\cN)$ は $\det$ の比で決まり、
$p\nmid n$ のとき Wakabayashi \cite[Prop.~4.1.4]{Wak} により $\cN$ の $n$ 乗根が
一意に取れるので、問題は $(\det\cF,\det\nabla)$ の持ち上げに帰着する。
ただし逆向き（$(\det\cF,\det\nabla)$ の持ち上げから $\mathrm{GL}_n$-oper 全体の
持ち上げを組み立てる部分）で $\det$ 以外に障害が無いことは完全には検証していない
（注意 \ref{rem:det-status}）。
\end{proof}

\begin{remark}[主張の地位と残るギャップ]\label{rem:det-status}
定理 \ref{thm:det} は $p\nmid n$ の仮定のもとで上のように証明される
\emph{[本研究・部分証明]}。系 \ref{cor:obstruction} も参照。
\GAP{「$\det$ 以外に持ち上げの障害が無い」ことの完全な議論が未了。
仮定を追加しても現状では閉じない（\texttt{handoff/level-chat.md} \S3.6, \S6.5）。}
\NeedsCheck{\cite[Prop.~4.1.4]{Wak} の正確な仮定と主張。}
なお $p\mid n$ の場合は本定理の仮定 $p\nmid n$ の対象外であり、
proof gap ではなく将来の拡張問題として \S\ref{sec:open} の問題 \ref{q:GP} で扱う。
\end{remark}

\begin{corollary}\label{cor:pi1}
定理 \ref{thm:det} のもとで、$\pi_1(G)$ が $p$ 群を含まない
（例えば $G$ 随伴型、あるいは $p\nmid n$ で $G=\mathrm{GL}_n$）場合、
レベル上昇の障害は $\Pic(X)/p^N$ の障害に帰着する。
特に $r\ge1$ ならば命題 \ref{prop:rank1curve} により障害は消える。
\end{corollary}

\begin{remark}[主張の地位]\label{rem:pi1-status}
系 \ref{cor:pi1} は定理 \ref{thm:det} に依存するので \emph{[本研究・部分証明]}。
\end{remark}

%%%%%%%%%%%%%%%%%%%%%%%%%%%%%%%%%%%%%%%%%%%%%%%%%%%%%%%%%%%%%%%%%%%%%%%%%%%%%%%
\section{反例族のリストと分類}\label{sec:list}

\begin{table}[htbp]
\centering
\caption{反例および非反例の一覧。「なし$^{\ast}$」は仮定 (W10.3)（\S\ref{sec:M03}）に
依存する条件付きの結論。}
\label{tab:examples}
\footnotesize
\setlength{\tabcolsep}{3.5pt}
\begin{tabular}{cllcl}
\toprule
\# & $(X,D)$・対象 & exponent/radius & 最大レベル & 障害の型\\
\midrule
1 & $(\bbP^1,\{0,\infty\})$; $(\cO,d+a\,\dlog x)$, $a\in\bbF_p^\times$
  & $a\in\bbZ/p$ & $\infty$ & なし（命題 \ref{prop:rank1curve}、部分証明）\\[2pt]
2 & 同上; $a\in k\setminus\bbF_p$ & --- & $0$ & Artin--Schreier \eqref{eq:AS}\\[2pt]
3 & $(\bbA^1,\{0\})$; $\nabla_\theta=\theta+N$, $N$ 冪零 $\neq0$
  & --- & $0$ & 留数の冪零成分（系 \ref{cor:nilpotent}）\\[2pt]
4 & $g\ge1$, $D=\emptyset$; $\cL=F^{(N)*}\cM$, $p\nmid\deg\cM$
  & $0$ に強制 & ちょうど $N$ & $\Pic$ の非 $p$ 可除性（例 \ref{ex:exactN}）\\[2pt]
5 & $(\bbP^1;0,1,\infty)$, $p=5$; 表 \ref{tab:p5} の $5$ 個
  & $\lambda=(3,3,3)$ 等 & $\infty^{\ast}$ & なし$^{\ast}$（定理 \ref{thm:M03}）\\[2pt]
6 & $X_{\mathbb{G}}$（totally deg.）; dormant $\mathrm{PGL}_2^{(N)}$-oper
  & $(s_e)\in\mathrm{Ed}_{p,N,\mathbb{G}}$ & $\infty^{\ast}$ & なし$^{\ast}$（系 \ref{cor:general}）\\[2pt]
7 & $g\ge1$, $D=\emptyset$; $\mathrm{GL}_2$-oper で $\det$ が \#4 型
  & --- & ちょうど $N$ & 行列式（定理 \ref{thm:det}、$p\nmid n$、部分証明）\\[2pt]
8 & $g\ge2$, $r=0$（一般 smooth）; dormant $\mathrm{PGL}_2$-oper
  & --- & \textbf{未決定} & \S\ref{sec:open} 問題 \ref{q:main}\\
\bottomrule
\end{tabular}
\end{table}

行 8 は当初「本命」と考えていたケースである（旧 ver.\,0.3 での問題 \ref{q:main}）。
行 6 と系 \ref{cor:general} により、totally degenerate な曲線の範囲では
（仮定 (W10.3) のもとで）反例が無いことが分かったが、
一般の smooth 曲線については未解決のままである。

\begin{proposition}[反例が存在しないための十分条件]\label{prop:sufficient}
次のいずれかが成り立つとき、dormant 対象は任意のレベルに持ち上がる。
\begin{enumerate}[label=(\alph*),leftmargin=2em]
\item 階数 $1$、$X$ 固有滑らかな曲線、$r\ge1$（命題 \ref{prop:rank1curve}）。
\item 留数が冪単：系 \ref{cor:nilpotent} により $\Res=0$ となり、
      非対数的な理論に帰着する（remark \ref{rem:res-zero}）。
\item $\mathrm{PGL}_2$-oper、$(g,r)=(0,3)$（定理 \ref{thm:M03}）。
\item $\mathrm{PGL}_2$-oper、$X$ が totally degenerate（系 \ref{cor:general}）。
\item 構造群 $G$ の $\pi_1(G)$ が $p$ 群を含まない（系 \ref{cor:pi1}）。
\end{enumerate}
ここで (a) は \emph{[部分証明]}（注意 \ref{rem:rank1curve-status}）、
(c), (d) は仮定 (W10.3)（\S\ref{sec:M03}）に依存する \emph{[条件付き]}、
(e) は定理 \ref{thm:det} に依存し、かつ $p\nmid n$ と「$\det$ 以外に障害が無いこと」
を要する \emph{[部分証明]}（注意 \ref{rem:det-status}）。
無条件なのは (b) のみである。
\end{proposition}

\begin{remark}[当初の予想の訂正]
本研究の出発点では
「対数格子の取り替え（$D$ での Hecke 変形）で回避できないことを示す」
ことが目標であった。
しかし補題 \ref{lem:hecke} と命題 \ref{prop:rank1curve} が示すとおり、
階数 $1$ ではむしろ\textbf{常に回避できる}。
当初は「反例の存在には格子の剛性（oper 構造）と $r=0$ の両方が必要」と考えたが、
$\mathrm{PGL}_2$-oper については $r=0$ でも（仮定 (W10.3) のもと、totally degenerate な
曲線上で）反例が存在しないことが分かった（系 \ref{cor:general}）。
訂正の経緯は付録 \ref{app:errata} を参照。
\end{remark}

%%%%%%%%%%%%%%%%%%%%%%%%%%%%%%%%%%%%%%%%%%%%%%%%%%%%%%%%%%%%%%%%%%%%%%%%%%%%%%%
\section{非可換 Hodge 対応との関係}\label{sec:NAHT}

命題 \ref{prop:padic} の構造は、
$\psi=0$ なる対数接続が $(\bbZ/p)^r$-次数付き $\cO_X$ 加群、
すなわち $\mu_p^r$-同変層として記述されることを意味する。
実際、$\nabla_{\theta_i}^{\,p}=\nabla_{\theta_i}$（補題 \ref{lem:toral}）から
固有分解 $\cE=\bigoplus_{a\in(\bbZ/p)^r}\cE_a$ が得られ、
$\nabla_{\theta_i}(x_ie)=x_i(\nabla_{\theta_i}+1)e$ より
$x_i$ は次数 $e_i$ をもつ。

\begin{conjecture}[root stack による Cartier 降下]\label{conj:rootstack}
$\mathcal{X}'_N:=\sqrt[p^N]{(X^{(N)},D^{(N)})}$（root stack）とおくと、
自然な log étale 射 $\widetilde F^{(N)}:X\to\mathcal{X}'_N$ が存在し、
\[
  \widetilde F^{(N)*}:\ \mathrm{QCoh}(\mathcal{X}'_N)
  \xrightarrow{\ \sim\ }\bigl\{\text{dormant レベル }N\text{ の }\Dlog{N-1}\text{-加群}\bigr\}
\]
が圏同値を与える。特に
$\widetilde F^{(N)*}\Omega^1_{\mathcal{X}'_N}(\log)\cong\Omega^1_X(\log D)$
であり（素朴な $F^*$ と異なり零にならない）、oper 構造は $\widetilde F^{(N)*}$ で保たれる。
\end{conjecture}

命題 \ref{prop:rank1} は予想 \ref{conj:rootstack} の階数 $1$・局所版であり、
$\cO(-\tilde a\,\mathcal{D})$（$\mathcal{D}$ は $p^N\mathcal{D}=D^{(N)}$ なる普遍因子）
の引き戻しが $(R,\nabla_a)$ に対応する。

この定式化は文献上、Lorenzon--Montagnon の
\emph{indexed algebra} $A^{gp}_X$, $\mathcal{B}^{(m+1)}_{X/S}$ として
既に与えられていると思われる。
Ohkawa \cite{Ohk} は、$X'\to S$ の $\bmod\,p^2$ 持ち上げの下で
$\cD^{(m)}$-作用つき indexed $A^{gp}_X$-加群と
Higgs 場つき indexed $\mathcal{B}^{(m+1)}_{X/S}$-加群の圏同値を構成しており、
これは Ogus--Vologodsky \cite{OV} と Schepler \cite{Sch} の結果のレベル $m$ 版である。
\Todo{予想 \ref{conj:rootstack} が Ohkawa の定理の言い換えに過ぎないのか、
それとも独立な主張なのかを確定する。}

%%%%%%%%%%%%%%%%%%%%%%%%%%%%%%%%%%%%%%%%%%%%%%%%%%%%%%%%%%%%%%%%%%%%%%%%%%%%%%%
\section{未解決問題}\label{sec:open}

\S\ref{sec:M03} により、当初の主要問題（問題 \ref{q:main}）は、
仮定 (W10.3) のもとで、かつ totally degenerate な曲線に限れば、
否定的に解決される。無条件の解決、および一般の smooth 曲線での解決は
未達である。

\begin{question}[旧・本命問題：条件付きで否定的に解決]\label{q:main}
種数 $g\ge2$ の非点付き曲線上で、
$F^1$-射影構造であって $F^2$-射影構造に持ち上がらないものを構成せよ
（当初の「本命」問題。旧 ver.\,0.3 では
「$F^N$-射影構造の個数公式を比較し $|{}^2\Op|<|{}^1\Op|$ となる $(g,p)$ を探す」
方針を掲げていた）。
系 \ref{cor:general} により、仮定 (W10.3) のもとでは、少なくとも
totally degenerate な曲線上にはそのような対象が存在しない。
残る空白は次の二点である：
(i) 仮定 (W10.3)（Wakabayashi \cite[Prop.~10.3.3]{Wak2} とレベル降下の両立性）
の確認；
(ii) 一般の smooth 曲線上での $\mathrm{res}_N$ の全射性
（注意 \ref{rem:no-pgl2-cex} の \Todo{}。generic étaleness \cite[Thm.~C]{Wak2}
からの直接の議論）。
\end{question}

\begin{question}[次数障害の十分性]\label{q:degree-suff}
命題 \ref{prop:degree} の合同式 \eqref{eq:degcong} は必要条件である。
$g\ge1$, $r\ge1$ において、これが（適当な安定性の下で）十分条件となるか。
$(g,r)=(0,3)$ では定理 \ref{thm:M03} が（仮定 (W10.3) のもとで）肯定的に答える。
\end{question}

\begin{question}[レベル $N$ の数え上げ公式]\label{q:count}
$\#\,{}^{\dagger}C_N$（$=\#\,{}^N\Op^{\Zzz}_{2,\sP}$、命題 \ref{prop:Cdagger}）の
閉じた表示を求めよ。$N=1$ では $p(p^2-1)/24$（命題 \ref{prop:count}）。
$N\ge2$ では \cite[Thm.~E]{Wak2} により $p$ の擬多項式 $H_{N,\mathbb{G}}(p)$
（$g=2$ のとき次数 $3N$）として存在することが分かっているが明示式はない。
計算済みの値（表 \ref{tab:count}、付録 \ref{app:num}）：
$p=3$ で $1,11,121$、$p=5$ で $5,225$、$p=7$ で $14,1666$。
$p=3$ の $1,11,121=11^2$ という並びは示唆的である。
条件 \ref{def:Cdagger}(b) の再帰構造から桁ごとの transfer matrix 表示が
得られないか。
\end{question}

\begin{question}[一般階数 $n$]\label{q:rank}
定理 \ref{thm:det} は任意の $n$ で（$p\nmid n$、および $\det$ 以外に障害が
無いことを措けば）成立するが、系 \ref{cor:general} の証明は $n=2$ の
generic étaleness \cite[Thm.~C]{Wak2} に依存している。
$n\ge3$ で $\Pi_{n,N,\rho,g,r}$ の generic étaleness を示せば、
$\mathrm{PGL}_n$ でも（仮定 (W10.3) のもとで）反例が存在しないことが従う。
これは \cite[\S1.6]{Wak2} の「canonical lifting は一般の $n$ でも存在すると
期待される」という問題と同値である。
\end{question}

\begin{question}[高次元]\label{q:higherdim}
$\dim X\ge2$、$D$ が SNC のとき exponent は $(\bbZ/p^N)^{\times r}$ に値をとる。
交差 $D_i\cap D_j$ における両立条件が新しい障害を生むか。
定理 \ref{thm:det} の類似（障害が $\det$ に局在する）は成り立つか。
曲線の場合と違い $\Pic$ の代わりに $H^2$ の捩れや Brauer 群が効く可能性がある。
Ohkawa の indexed algebra の枠組みで定式化できると思われる。
\end{question}

\begin{question}[一般 $(G,P)$]\label{q:GP}
radius は $\mu_{p^N}\to T$ の $W$-共役類として定義される。
系 \ref{cor:pi1} を精密化し、
「レベル上昇の障害 $=$ $H^1(X,\pi_1(G))$ の $p^N$ 可除性の障害」
という形の一般命題を定式化・証明せよ。
$\mathrm{PGL}_n$ では $\pi_1=\bbZ/n$ であり、$\gcd(n,p)=1$ なら障害は消える
（\cite[Prop.~4.1.4]{Wak} が $p\nmid l$ を仮定していることに対応）。
$p\mid n$ の場合が未検討であり、これは定理 \ref{thm:det} の仮定 $p\nmid n$ を
除去する問題でもある。
\end{question}

\begin{question}[root stack による定式化]\label{q:rootstack}
予想 \ref{conj:rootstack} が Ohkawa \cite{Ohk} の定理の言い換えに過ぎないのか、
独立な主張なのかを確定せよ。確定すれば、\S\ref{sec:local} の exponent の
$p$ 進展開が $\mathcal{X}'_N=\sqrt[p^N]{(X^{(N)},D^{(N)})}$ 上の
$\mu_{p^N}$ 次数付けとして幾何的に説明される。
\end{question}

\begin{question}[中間領域]\label{q:middle}
系 \ref{cor:nilpotent} により、dormant な世界（留数 toral）と
nilpotent indigenous bundle の世界（留数冪単）は交わらない。
$p^N$-曲率が冪零だが非零であるクラスに、両者を繋ぐ構造があるか。
\end{question}

%%%%%%%%%%%%%%%%%%%%%%%%%%%%%%%%%%%%%%%%%%%%%%%%%%%%%%%%%%%%%%%%%%%%%%%%%%%%%%%
\appendix
\section{規約の対照表}

\begin{table}[htbp]
\centering
\begin{tabular}{lll}
\toprule
& Berthelot & Wakabayashi \cite{Wak}（本稿）\\
\midrule
作用素環 & $\cD^{(m)}$ & 「レベル $N$」$=\cD^{(N-1)}$\\
曲率 & $p^{m+1}$-曲率 $\nabla_{\langle p^{m+1}\rangle}$ & 「$p^N$-曲率」\\
dormant & $p^{m+1}$-曲率 $=0$ & 「dormant レベル $N$」\\
非対数での意味 & $F^{(m+1)*}$ の像 & $F^{(N)*}$ の像\\
radius の値域 & --- & $(\bbZ/p^N\bbZ)^\times/\{\pm1\}$（$n=2$）\\
\bottomrule
\end{tabular}
\end{table}

正標数 $\bbZ_p$ 値留数の研究で用いられる規約（$\cD_U^{(n)}$ を位数 $<p^n$ で定義）
および logarithmic holonomy の研究で用いられる規約
（$\cD^{(m)}_{(X,D)}:=\mathcal{E}nd_{\cO_{X^{(m+1)}}}(\cO_X)$）との対照は
\texttt{notes/topics/00-common/conventions.tex} を参照。本稿ではそれらを採用・統一しない
（注意 \ref{rem:conv-status}）。

\section{訂正の履歴}\label{app:errata}

研究の途上で誤りと判明した主張・方針を、研究史として残す。
出典は \texttt{handoff/level-chat.md}（特に \S5 と付録B）および
\texttt{handoff/app-session-2026-09-03.md} \S8.1。

\begin{enumerate}[leftmargin=1.8em]
\item \textbf{「log では『$\cD^{(1)}$-作用の存在』と『$p$-曲率 $0$』が乖離する」（誤り）.}\
      素朴な Frobenius 引き戻し $F^*$ と $\cD^{(1)}$-加群構造への持ち上げを混同していた。
      $\bbZ_p$ の恒等式 $\prod_{j=0}^{p-1}(\theta-j)=p!\binom{\theta}{p}$ の左辺が
      $\bmod\,p$ で $\theta^p-\theta$ となるため、$a\in\bbF_p^\times$ の対象も
      $\cD^{(1)}$-加群として持ち上がる（注意 \ref{rem:binom}, \ref{rem:early-misread}）。
      規約のズレであって数学的現象ではなかった。
\item \textbf{「$D$ での Hecke 変形で回避できないことを示す」という当初のタスク（逆だった）.}\
      補題 \ref{lem:hecke} と命題 \ref{prop:rank1curve} の通り、階数 $1$・$r\ge1$ では
      格子の取り替えで exponent を自由に動かせるので、局所的な障害は常に回避できる。
      対数構造（marked point）は障害を消す方向に働く。
\item \textbf{$\mathfrak{sl}_2$ level $k=p-2$ の WZW 融合則による数え上げ（誤り）.}\
      $\binom{p+1}{3}$（$p=5$ で $20$）と主張したが正しくは $p(p^2-1)/24$（$p=5$ で $5$）。
      radius は $\bbF_p^\times/\{\pm1\}$（$(p-1)/2$ 個）に値をとるのに
      $\{0,\dots,p-2\}$（$p-1$ 個）の可積分ウェイトと対応させた誤り。
      本稿は WZW を用いない。
\item \textbf{$N=2$ の数え上げ $610$（$p=5$）, $2262$（$p=7$）（誤り、かつ未計算）.}\
      前稿 ver.\,0.3 の表 \ref{tab:count} にこれらを「計算機による直接列挙」として
      記載したが、列挙は実行されておらず値も誤り。正しい値は $225$, $1666$
      （注意 \ref{rem:erratum}、付録 \ref{app:num}）。素朴な $\lambda$ 定式化による
      過大評価値 $345$, $3122$ とも一致しない。
\item \textbf{数値を「計算機で確認」と称して未実行だったこと.}\
      ver.\,0.4 では、確定値として記す数値はすべて付録 \ref{app:num} のスクリプトを
      実行した結果に限る。
\item \textbf{当初の主要予想「$r=0$, $g\ge2$ に $F^1$ から $F^2$ に持ち上がらない
      $\mathrm{PGL}_2$-oper がある」.}\
      系 \ref{cor:general}（仮定 (W10.3) のもと、totally degenerate な曲線）により
      否定的な答えが強く示唆される。ただし無条件ではなく、一般の smooth 曲線でも
      未解決（問題 \ref{q:main}）。
\end{enumerate}

\section{文献確認事項（未了）}\label{app:litcheck}

本稿の主張が依存する、まだ原論文で確認していない事項。優先度順。

\begin{enumerate}[leftmargin=1.8em]
\item \textbf{[最優先, L2]} Wakabayashi \cite{Wak2}（arXiv:2209.08528）\S10.3、
      特に Prop.~10.3.3 の径数付けとレベル降下の両立性。
      $\Rightarrow$ 仮定 (W10.3)、したがって定理 \ref{thm:M03}, 系 \ref{cor:general},
      命題 \ref{prop:sufficient}(c),(d)。確認できれば定理 \ref{thm:M03} が
      無条件になる（\texttt{handoff/level-chat.md} \S3.3, \S7 \#2）。
\item \textbf{[書誌]} arXiv:2209.08528 の題名。ar5iv 上では
      ``Canonical diagonal liftings I: dormant $\mathrm{PGL}_n$-opers on algebraic curves''
      であり、本稿の \cite{Wak2} の題名と一致しない。改題か引用誤りか未判定
      （\texttt{handoff/app-session-2026-09-03.md} \S8.3）。
\item \cite{Wak2} の Prop.~10.1.4（命題 \ref{prop:Cdagger}）、Thm.~E
      （注意 \ref{rem:genus2}、系 \ref{cor:general}）、(1.2)、\S1.6、Thm.~C
      （generic étaleness、注意 \ref{rem:no-pgl2-cex}、問題 \ref{q:rank}）。
\item Berthelot \cite[\S1.1--1.2]{Ber1}：$\cD^{(m)}$ の定義
      （$\partial^{[j]}\in\cD^{(m)}$ の $j$ の範囲）。$\Rightarrow$ 規約の基準
      （注意 \ref{rem:conv-status}）。
\item Wakabayashi \cite{Wak}（arXiv:2201.11266）の Prop.~2.1.5, Cor.~2.1.6,
      Def.~2.1.3, Prop.~3.2.1, Rem.~3.3.3, Prop.~3.3.4, Prop.~3.4.1, Prop.~4.1.4,
      Def.~4.2.1, Def.~4.3.2, Thm.~C。\S\ref{sec:local}, \S\ref{sec:opers} は
      これらの引用に依拠しており、ar5iv 経由の二次確認にとどまる
      （\texttt{handoff/app-session-2026-09-03.md} \S0.3）。
\item Mochizuki \cite{Moc}, Chap.~II, Prop.~1.4, 1.5（命題 \ref{prop:radius-unit}）。
      Wakabayashi 経由の引用のみ。
\item Montagnon \cite{Mon} の thèse における対数版 $\cD^{(m)}$ の presentation
      （注意 \ref{rem:binom} の $\binom{\theta}{p}$ 議論の裏付け）。
\item Ohkawa \cite{Ohk}：予想 \ref{conj:rootstack} が Ohkawa の定理の
      言い換えに過ぎないか（問題 \ref{q:rootstack}）。
\item Osserman \cite{Oss} の正確な巻号（\S\ref{app:num} には無関係、書誌のみ）。
\item \texttt{check.py}（\texttt{handoff/level-chat.md} 付録Aが言及）は \texttt{handoff/}
      には含まれていなかったが、作業環境の \texttt{\~{}/Downloads/check.py}
      に原ファイルが残っており、\texttt{computations/check.py} として回収した
      （\S\ref{app:num}）。
\end{enumerate}

\section{数値検証}\label{app:num}

本稿の数値は次のスクリプトを本稿の作業環境（Python 3）で実行して得た。

\begin{itemize}[leftmargin=1.8em]
\item \texttt{computations/count\_dagger\_C.py}：定義 \ref{def:Cdagger} から
      独立に書き起こした ${}^{\dagger}C_N$ の列挙器。
      \texttt{in\_dagger\_C(s,N,p)} が条件 (a),(b) をそのまま判定し、
      \texttt{count\_dagger\_C(N,p)} が $\{0,\dots,p^N-1\}^3$ を総当たりする
      （条件 (a) より $s_i\le p^N-2$ なのでこれで尽くされる）。
\item \texttt{computations/check.py}：\texttt{handoff/level-chat.md} 付録Aが言及する
      原スクリプトを作業環境の \texttt{\~{}/Downloads/} から回収したもの
      （改変なし）。
\item \texttt{computations/verify\_p5.py}（$=$ \texttt{handoff/verify\_p5.py}）：
      表 \ref{tab:p5} の $5$ 個の被覆について、$\bbF_5$ 上で Wronskian と
      各 ODE が $f,g$ を解にもつことを検証する。
\end{itemize}

\noindent 実行結果（両列挙器で一致）：

\begin{center}
\begin{tabular}{lrrrrr}
\toprule
$p$ & 3 & 5 & 7 & 11 & 13\\
\midrule
$\#\,{}^{\dagger}C_1$ & 1 & 5 & 14 & 55 & 91\\
$p(p^2-1)/24$ & 1 & 5 & 14 & 55 & 91\\
$\#\,{}^{\dagger}C_2$ & 11 & 225 & 1666 & 24805 & 67431\\
素朴な $\lambda$（条件 (b) なし）& 11 & 345 & 3122 & 55605 & 158431\\
\bottomrule
\end{tabular}
\end{center}

\noindent さらに $\#\,{}^{\dagger}C_3\big|_{p=3}=121=11^2$。
入れ子性 ${}^{\dagger}C_N\subseteq{}^{\dagger}C_{N+1}$ は
$(p,N)\in\{(3,1),(5,1),(3,2),(5,2)\}$ で
${}^{\dagger}C_N\setminus{}^{\dagger}C_{N+1}=\emptyset$ を確認。
$p=5$, $N=1$ の $(s_1,s_2,s_3)$ 全リストは
$(0,0,0),(0,1,1),(1,0,1),(1,1,0),(1,1,1)$ で、
$\lambda_i=2s_i+1$ とすると表 \ref{tab:p5} の $5$ 個に一致する。
\texttt{verify\_p5.py} は $5$ 個すべてについて $f,g$ が対応 ODE の解であること、
Wronskian が \eqref{eq:wronskian} と整合することを確認した。

\medskip
\noindent これらは一般 $p$（あるいは一般 $N$）に対する証明ではなく、
上記の有限個の $(p,N)$ での計算機確認である。

%%%%%%%%%%%%%%%%%%%%%%%%%%%%%%%%%%%%%%%%%%%%%%%%%%%%%%%%%%%%%%%%%%%%%%%%%%%%%%%
\begin{thebibliography}{99}

\bibitem{Ber1} P.~Berthelot,
\emph{$\mathcal{D}$-modules arithmétiques.\ I.\ Opérateurs différentiels de niveau fini},
Ann.\ Sci.\ École Norm.\ Sup.\ (4) \textbf{29} (1996), 185--272.

\bibitem{Ber2} P.~Berthelot,
\emph{$\mathcal{D}$-modules arithmétiques.\ II.\ Descente par Frobenius},
Mém.\ Soc.\ Math.\ Fr.\ (N.S.) \textbf{81} (2000).

\bibitem{Kat} N.~Katz,
\emph{Nilpotent connections and the monodromy theorem: applications of a result of Turrittin},
Publ.\ Math.\ IHÉS \textbf{39} (1970), 175--232.

\bibitem{Moc} S.~Mochizuki,
\emph{Foundations of $p$-adic Teichmüller theory},
AMS/IP Studies in Advanced Mathematics \textbf{11}, AMS, 1999.

\bibitem{Mon} C.~Montagnon,
\emph{Généralisation de la théorie arithmétique des $\mathcal{D}$-modules à la géométrie logarithmique},
Thèse, Université de Rennes I, 2002.

\bibitem{Ohk} S.~Ohkawa,
\emph{On logarithmic nonabelian Hodge theory of higher level in characteristic $p$},
Rend.\ Sem.\ Mat.\ Univ.\ Padova \textbf{134} (2015), 47--91.

\bibitem{Ohk2} S.~Ohkawa,
\emph{On log local Cartier transform of higher level in characteristic $p$},
arXiv:1410.0535.

\bibitem{OV} A.~Ogus and V.~Vologodsky,
\emph{Nonabelian Hodge theory in characteristic $p$},
Publ.\ Math.\ IHÉS \textbf{106} (2007), 1--138.

\bibitem{Oss} B.~Osserman,
\emph{Frobenius-unstable bundles and $p$-curvature},
Trans.\ Amer.\ Math.\ Soc.\ \textbf{360} (2008), 273--305.
\NeedsCheck{正確な巻・年・頁を確認（付録 \ref{app:litcheck}）}

\bibitem{Sch} D.~Schepler,
\emph{Logarithmic nonabelian Hodge theory in characteristic $p$},
arXiv:0802.1977.

\bibitem{Wak} Y.~Wakabayashi,
\emph{Differential modules and dormant opers of higher level},
arXiv:2201.11266.

\bibitem{Wak2} Y.~Wakabayashi,
\emph{Arithmetic liftings and 2d TQFT for dormant opers of higher level},
arXiv:2209.08528.
\NeedsCheck{ar5iv 上の題名は
``Canonical diagonal liftings I: dormant $\mathrm{PGL}_n$-opers on algebraic curves''
であり不一致。改題か引用誤りか未判定。\S10.1--10.6, \S9.2 は未読
（付録 \ref{app:litcheck}）。}

\bibitem{Wak3} Y.~Wakabayashi,
\emph{A theory of dormant opers on pointed stable curves},
Astérisque \textbf{432} (2022).
\NeedsCheck{radius の定義（Chap.~2, \S2.8, Def.~2.32 など）は未確認。}

\end{thebibliography}

\end{document}
